% Part: methods
% Chapter: proofs
% Section: cant-do-it

\documentclass[../../../include/open-logic-section]{subfiles}

\begin{document}

\olfileid{mth}{prf}{cnt}

\olsection{از پسش برنمی‌آیم!{}}

همهٔ ما به نقطه‌ای می‌رسیم که احساس می‌کنیم باید تسلیم شویم. اما شما
\emph{می‌توانید} از پسش برآیید. مدرس، دستیار آموزشی و هم‌کلاسی‌هایتان
می‌توانند کمک کنند. از آن‌ها کمک بخواهید!{} در ادامه چند پیشنهاد آمده
است تا از رسیدن به بحران جلوگیری کنید و اگر احساس کردید می‌خواهید
تسلیم شوید، بدانید چه کنید.

برای آنکه بتوانید مسئله‌ها را با موفقیت حل کنید، کارهای زیر را انجام دهید:
\begin{enumerate}
\item \emph{تا حد ممکن زود آغاز کنید.} در طول نیم‌سال سرمان شلوغ
  می‌شود و بسیاری از ما با به‌تعویق‌انداختن کارها دست‌وپنجه نرم
  می‌کنیم. یکی از بهترین کارهایی که می‌توانید انجام دهید این است که
  تکلیف‌ها را زود آغاز کنید. در این صورت اگر گیر کردید، برای یافتن راه
  حل وقت دارید و لازم نیست به گریه پناه ببرید.
\item \emph{با هم‌کلاسی‌هایتان گفت‌وگو کنید.} تنها نیستید. دیگران در
  کلاس نیز ممکن است مشکل داشته باشند، هرچند شاید مشکلشان با چیزهای
  دیگری باشد. گفت‌وگو با هم‌کلاسی‌ها می‌تواند چشم‌اندازی متفاوت دربارهٔ
  مسئله به شما بدهد و به گره‌گشایی بینجامد. البته پاسخشان را صرفاً
  رونویسی نکنید: از آن‌ها راهنمایی بخواهید، یا توضیح دهید کجا گیر
  کرده‌اید و گام بعدی را بپرسید. وقتی هم راه را یافتید، کمک را جبران
  کنید. همراهی با شخصی دیگر و توضیح‌دادن مطالب، فهم خودتان را نیز بهتر
  می‌کند.
\item \emph{کمک بخواهید.} منابع بسیاری در اختیار دارید؛ مدرس و دستیار
  آموزشی برای کمک به شما آنجا هستند و \emph{می‌خواهند} موفق شوید.
  آن‌ها باید بتوانند در حل مسئله همراهی‌تان کنند و تشخیص دهند در کدام
  بخش فرایند مشکل دارید.
\item \emph{استراحت کنید.} اگر گیر کرده‌اید، \emph{ممکن است} علتش این
  باشد که بیش از اندازه به مسئله خیره مانده‌اید. کمی استراحت کنید،
  فنجانی چای بنوشید، یا مدتی روی مسئله‌ای دیگر کار کنید؛ سپس با ذهنی
  تازه به مسئله بازگردید. بگذارید یک شب بگذرد.
\end{enumerate}

توجه کردید این راهبردها مستلزم آن‌اند که کار روی اثبات را از مدت‌ها
پیش آغاز کرده باشید؟ اگر ساعت دو بامدادِ روز پیش از موعد تحویل شروع
کرده باشید، شاید این پیشنهادها چندان سودمند نباشند.

ممکن است این سخنان بسیار بدبینانه به نظر برسد، اما حل یک اثبات چالشی
است که در پایان ارزشش را نشان می‌دهد. برخی افراد این کار را به‌عنوان
حرفه انجام می‌دهند؛ پس حتماً جنبه‌ای لذت‌بخش هم دارد. مانند تقریباً هر
کار دیگر، حل مسئله و انجام اثبات نیازمند تمرین است. شاید هم‌کلاسی‌هایی
را ببینید که این کار برایشان آسان است؛ احتمالاً فقط از پیش بسیار تمرین
کرده‌اند. خیلی زود تسلیم نشوید.

اگر برای مسئله‌ای خاص وقت یا حوصله‌تان تمام شد، اشکالی ندارد. این به
آن معنا نیست که نادانید یا هرگز آن را نخواهید فهمید. از مدرس یا
دانشجویی دیگر بپرسید چگونه حل می‌شود و مشخص کنید کجا اشتباه کردید یا
گیر افتادید، تا دفعهٔ بعد که با مسئله‌ای مشابه روبه‌رو شدید از آن
پرهیز کنید. سپس بکوشید بدون نگاه‌کردن به پاسخ خودتان حلش کنید. دفعهٔ
بعد نیز زودتر آغاز کنید و زودتر کمک بخواهید.

\end{document}
